\documentclass[11pt]{article}

\usepackage[margin=1in]{geometry}
\usepackage{booktabs}
\usepackage{longtable}
\usepackage{array}
\usepackage{xcolor}
\usepackage{hyperref}
\usepackage{enumitem}

\hypersetup{
  colorlinks=true,
  linkcolor=blue,
  urlcolor=blue
}

\setlist[itemize]{leftmargin=1.2em}
\setlist[enumerate]{leftmargin=1.4em}
\renewcommand{\arraystretch}{1.18}

\newcolumntype{P}[1]{>{\raggedright\arraybackslash}p{#1}}
\newcommand{\status}[1]{\textbf{#1}}

\title{FCFS Appointment Simulation\\Final Hypothesis Synthesis Summary}
\author{CUIMC Appointment Simulation}
\date{May 29, 2026}

\begin{document}
\maketitle

\section*{Purpose}

This document summarizes the final professor-facing hypothesis synthesis for
the FCFS appointment simulation. Earlier reports were investigation notes. The
final synthesis states what the implemented simulation proves or disproves about
access, utilization, delay, no-offer, cancellation, balking, no-show, and class
gaps.

The final deliverable is:

\begin{itemize}
  \item \texttt{hypothesis\_synthesis\_professor\_comparison.html}
  \item source: \texttt{hypothesis\_synthesis\_professor\_comparison.qmd}
  \item audit record: \texttt{supporting/hypothesis\_synthesis\_professor\_comparison\_audit.md}
\end{itemize}

\section*{Final Audit Verdict}

The final synthesis passes audit for the current implemented model. Three
independent checks were run:

\begin{enumerate}
  \item evidence and numeric consistency against conclusion reports and output
        CSV files;
  \item implementation consistency against the simulation code;
  \item structure and redundancy review for a professor-facing conclusion
        format.
\end{enumerate}

No material numeric mismatches were found. No additional experiments are needed
to support the final nine hypotheses within the current model. Further work is
only needed for external validity or policy extensions, such as clinic
calibration, alternative no-show timing, or non-FCFS priority rules.

\section*{Minimum Reporting Set}

The simulation should not be summarized with utilization alone. The minimum
reporting set is:

\begin{itemize}
  \item served rate;
  \item utilization;
  \item offered wait;
  \item no-offer share;
  \item outcome decomposition;
  \item class-level served rates.
\end{itemize}

At the stress-test baseline, the outcome decomposition is approximately:

\begin{center}
\begin{tabular}{lr}
\toprule
Outcome & Share of arrivals \\
\midrule
Served & 26.9\% \\
Balked & 35.5--35.6\% \\
Canceled & 32.3--32.5\% \\
No-show & 5.1\% \\
No-offer & 0.0\% \\
\bottomrule
\end{tabular}
\end{center}

This means baseline losses are mainly post-offer losses, not horizon-denial
losses.

\section*{Implementation Mechanics}

\begin{longtable}{P{0.34\textwidth} P{0.58\textwidth}}
\toprule
Implemented mechanic & Consequence for interpretation \\
\midrule
\endfirsthead
\toprule
Implemented mechanic & Consequence for interpretation \\
\midrule
\endhead
Daily arrivals are generated by class, randomly permuted, and processed through
one pooled day-level FCFS queue. &
Class labels have no priority by construction. Class gaps require behavioral
asymmetry or interactions. \\

The engine offers the earliest residual day with available capacity, including
same-day slots. &
Offered wait is a congestion metric conditional on receiving an offer. \\

If the finite horizon is full, the current and remaining daily arrivals are
counted as no-offer. &
No-offer is a distinct access failure and must be reported with wait metrics. \\

Offered delay is recorded before balking; accepted delay is recorded only after
acceptance. &
Accepted wait can improve by selection when high-delay patients balk. \\

Balking happens before booking and consumes no slot. &
Balking reallocates access under high demand; it does not create capacity. \\

Cancellations apply to future bookings before new arrivals are processed. &
Future cancellations can free rebookable capacity for later arrivals. \\

No-shows are resolved on the service day and are not rebooked. &
No-show slots are unrebookable completed-visit losses. \\

No-show probability depends on original offered delay, not remaining wait. &
Demand level and threshold placement determine when no-show effects become
material. \\
\bottomrule
\end{longtable}

\section*{Final Hypotheses}

\begin{longtable}{P{0.07\textwidth} P{0.44\textwidth} P{0.13\textwidth} P{0.28\textwidth}}
\toprule
ID & Final conclusion & Verdict & Evidence basis \\
\midrule
\endfirsthead
\toprule
ID & Final conclusion & Verdict & Evidence basis \\
\midrule
\endhead
P1 &
Under oversubscription, utilization can remain high while patient access is
poor; served rate is the primary aggregate access metric. &
\status{True} &
Baseline utilization is about 0.841 while served rate is about 0.269. With
100 arrivals/day and 32 slots/day, served rate is mechanically capped near
0.32 before behavioral losses. \\

P2 &
With fixed capacity, increasing arrivals lowers served rate and raises offered
wait; the stress-test baseline is for mechanism diagnosis, not clinic
calibration. &
\status{True, scoped} &
The research note reports arrival rate as the dominant driver of served rate
($\beta=-0.774$) and offered wait ($\beta=+0.576$). \\

P3 &
Offered wait is a better congestion diagnostic than accepted wait, but it must
be reported with no-offer share. &
\status{True} &
Offered delay is recorded before balking; no-offer patients have no delay
observation. No-balking diagnostic no-offer rises to 0.2966. \\

P4 &
Balking reduces affected-class access, improves accepted-wait statistics by
selection, and reallocates capacity; it does not recover aggregate capacity
under the stress-test baseline. &
\status{True} &
Class 1 balking sweep: accepted delay falls from 9.95 to 6.71 days and Class 1
served rate falls from 0.297 to 0.237, while aggregate served rate stays near
0.269. \\

P5 &
Under high demand, future cancellations are reabsorbed by the pool: the
canceling class loses access while other patients may gain nearer slots. &
\status{Qualified true} &
Cancellation sweep: offered delay falls from 11.84 to 7.47 days; Class 1 served
rate falls from 0.344 to 0.175; Class 2 served rate rises from 0.135 to 0.395. \\

P6 &
No-shows are unrebookable booked-slot losses. In the tested regimes, utilization
effects become material when offered delays frequently cross the no-show-risk
region or when low-delay no-show probability is nonzero. &
\status{Qualified true} &
Stress-test no-show sweep drops utilization from 0.920 to 0.681. At realistic
$\lambda/S=1.2$, doubling no-show intensity changes utilization only from
0.935 to 0.934. \\

P7 &
Within the tested symmetric arrival-share range, pooled day-level FCFS does not
create a meaningful per-arrival class access gap from class labels or arrival
share alone. &
\status{True} &
The symmetric arrival-share experiment keeps the gap between about -0.002 and
+0.002 as Class 1 share varies from 0.30 to 0.70. \\

P8 &
Asymmetric behavior can create large class gaps, but driver ranking is
regime-specific. The broad regression screen emphasizes cancellation gap; the
realistic local attribution emphasizes no-show asymmetry. &
\status{Qualified true} &
Broad screen: cancellation gap has the largest absolute class-gap coefficient
(0.452). Realistic attribution: no-show contributes 0.120 of the 0.213 gap,
cancellation 0.073, balking threshold 0.019. \\

P9 &
Booking horizon length can amplify the realistic class gap until roughly a
three-week horizon, then the effect plateaus. &
\status{True} &
Horizon sweep: class gap rises from 0.167 at 7 days to 0.203 at 14 days and
0.213 at 21 days, then remains 0.213 at 28 and 42 days. \\
\bottomrule
\end{longtable}

\section*{Initial Hypotheses vs Final Hypotheses}

\begin{longtable}{P{0.30\textwidth} P{0.31\textwidth} P{0.31\textwidth}}
\toprule
Initial hypothesis or draft row & Final hypothesis or conclusion & Why it changed \\
\midrule
\endfirsthead
\toprule
Initial hypothesis or draft row & Final hypothesis or conclusion & Why it changed \\
\midrule
\endhead
H1: A class with higher cancellation probability shortens offered delay for
everyone but lowers its own served rate. &
P5: Under high demand, future cancellations are reabsorbed by the pool: the
canceling class loses access while other patients may gain nearer slots. &
The evidence supports the mechanism, but only with a high-demand rebooking
qualifier. It should not be generalized as ``cancellations are good.'' \\

H2: No-show events are pure loss: they reduce utilization without improving any
other metric. &
P6: No-shows are unrebookable booked-slot losses, but their material utilization
effect depends on delay exposure and low-delay no-show probability. &
The stress-test sweep supports no-show as unrebookable loss, but the realistic
demand grid shows almost no utilization movement at $\lambda/S=1.2$. \\

H3: Higher balking decreases accepted delay but decreases served rate; Class 2
accepted delay may rise at low-to-moderate Class 1 balking. &
P4: Balking reduces affected-class access, improves accepted-wait statistics by
selection, and reallocates capacity rather than recovering aggregate capacity. &
Accepted wait improves because long-delay patients leave the accepted sample.
Aggregate served rate can stay nearly flat when the other class backfills. \\

H4: Balking losses preserve utilization; no-show losses reduce utilization;
served rate decreases in both cases. &
P4, P5, P6: Loss timing is the real mechanism: balking is pre-booking,
cancellation is future-slot rebookable, and no-show is service-day
unrebookable. &
The utilization contrast is correct, but the served-rate claim needs an
affected-class qualifier. H4 was a synthesis of three mechanisms, not one
independent final hypothesis. \\

WIP H5: Cancellation reopens future capacity and can shorten delays. &
P5: Future cancellations can free rebookable capacity under high demand, but
they reduce the canceling class's served rate. &
This was redundant with H1. The final version keeps one cancellation row with
the class-impact and regime qualifiers. \\

WIP H6: A class with more arrivals may occupy more slots but also has more
arrivals in the denominator. &
P2 and P7: Total arrival pressure lowers aggregate access, but arrival share
alone does not create a meaningful per-arrival class gap under symmetric
behavior. &
The broad capacity-pressure result is true for total demand. The controlled
arrival-share experiment refutes arrival share alone as a class-gap mechanism. \\

WIP H7: Symmetric changes mostly affect aggregate outcomes; class advantage
emerges when one class has more favorable behavior. &
P7 and P8: Pooled FCFS is neutral under symmetric behavior; asymmetric behavior
can create large class gaps, and the dominant driver is regime-specific. &
The final version separates two questions: whether pooled FCFS creates gaps
without asymmetry, and which asymmetry drives the gap in each regime. \\

No explicit initial hypothesis on utilization vs access. &
P1: Under oversubscription, utilization can remain high while patient access is
poor; served rate is the primary aggregate access metric. &
This emerged as the central conclusion from the conclusion package and research
note. It is a denominator issue: utilization is slot-based, served rate is
arrival-based. \\

No explicit initial hypothesis on no-offer reporting. &
P3: Offered wait must be reported with no-offer share because no-offer patients
are excluded from wait metrics. &
The code and no-balking diagnostic show that wait metrics are conditional on
receiving an offer. No-offer is therefore a distinct access failure. \\

No explicit initial hypothesis on booking horizon. &
P9: Booking horizon length can amplify the realistic class gap until roughly a
three-week horizon, then the effect plateaus. &
This came from the newer horizon sweep. It is a design lever, not just another
behavioral asymmetry. \\
\bottomrule
\end{longtable}

\section*{Key Caveats}

\begin{itemize}
  \item The stress-test baseline is for mechanism diagnosis, not a calibrated
        clinic estimate.
  \item Regression coefficients are standardized associations over sampled
        parameter configurations, not standalone policy causal effects.
  \item Cancellation is a repeated daily future-booking hazard, not a one-time
        per-booking probability.
  \item No-show depends on original offered delay in the implemented model.
  \item Same-day cancellations are not modeled.
  \item Any policy claim about priority scheduling, overbooking, or calibrated
        clinic operation would require additional experiments.
\end{itemize}

\section*{Bottom Line}

The final conclusion is that pooled FCFS is mechanically neutral under symmetric
behavior, but aggregate metrics can hide access failures and class gaps when
behavior differs by class. Utilization, served rate, offered wait, no-offer,
outcome decomposition, and class-level served rates must be reported together.
Balking, cancellation, and no-show cannot be treated as interchangeable losses:
they occur at different points in the queueing pipeline and therefore affect
capacity, access, and class gaps differently.

\end{document}
